
De plus, on a une structure [$\pi$] privilégiée sur $G^{0}_{\mathbb{R}}$,
qui peut \struck{définir} s'exprimer définie par un hom
\struck{\ill{}}
\[
  i : \mu_{2,\mathbb{C}} \longrightarrow G^{0}_{\mathbb{C}}
\]
(comme on \ill{} en envoyant un tore [\ill{}] (\ill{})). D'autre part,
$G(s')_{\mathbb{C}}$ est muni de $i(s') : \mathbb{G}_{m,\mathbb{C}} \to
G(s')_{\mathbb{C}}$, et par restriction [à $\mu_2$] on obtient
canoniquement
\[
  \psi(s') = u(s')_{\mathbb{C}} \circ \bigl(i(s') \vert \mu_2\bigr) :
  \mu_{2,\mathbb{C}} \longrightarrow G_{\mathbb{C}}^{\xi(s')}
\]
\struck{D'autre part} Ceci posé, il faut que $\pi^{\xi(s')}$ (\ill{} de
$G_{\mathbb{R}}^{\xi(s')}$ déduite de $\pi$ sur $G$ par transport) soit
celle définie par les $\mu_{2,\mathbb{C}}$.
\note{la lettre du premier membre, $\psi(s')$, est douteuse : une sorte de
$\mathrm{v}$ à queue}

Par conséquent, puisqu'il s'agit [de \ill{}] de $G_{\mathbb{R}}^{\xi(s')}$
[\ill{}] que $\xi(s')$ \ill{} [\uncertain{à conjugaison près}], les
$\pi$-structures en question doivent être identiques --- \ill{} $u(s')$
doivent être une condition non tautologique [\ill{}] ce qui donne une
[\uncertain{restriction}] sur $u(s')$). [Notons cependant que la
distribution des $\pi$-structures d'un gpe réel $H$, définie en termes des
hom.\ $\mu_{2,\mathbb{C}} \to H_{\mathbb{C}}$, \ill{} \ill{}
\uncertain{commutatif} en tous cas !]
\note{« identiques » et « tautologique » portent un trait au-dessus ; dans
la parenthèse finale, un mot est biffé avant « $\pi$-structures »}

? d'\uncertain{ailleurs} \ill{} \uncertain{finie}, [si \ill{}] \ill{}
\ill{} \ill{} \uncertain{dominance}. \ill{} les \uncertain{restrictions}
\ill{} \ill{}
\note{deux lignes rapides, précédées d'un « ? » dans la marge gauche, dont
presque rien ne se lit}

Enfin, conjecturalement \ill{} [\uncertain{avec les conjectures} de]
Tate vérifiées, la connaissance des $u(s')$ \struck{\ill{}} \ill{} de la
structure d'effectivité dans la catégorie

